\subsection{Exact finite targets under independent digital shifts}
\label{sec:rqmc-finite-target}

I fix a certified financial graph before generating any random shifts. Its input
group is $\mathcal G=\prod_{j=1}^{d}\{0,\ldots,2^{b_j}-1\}$, with addition given
by coordinatewise bitwise exclusive-or, denoted by $\oplus$. The retained graph
returns an integer $a(u)$, and I write $x(u)=2^{-q}a(u)$. The finite range
certificate gives $0\leq x(u)\leq\rho$ for every $u\in\mathcal G$. I denote
its exact uniform mean by
\[
 \mu=|\mathcal G|^{-1}\sum_{u\in\mathcal G}x(u).
\]
For any deterministic points $t_1,\ldots,t_N\in\mathcal G$, a shift
$\Delta$ that is exactly uniform on $\mathcal G$ defines
\[
 B(\Delta)=\frac1N\sum_{i=1}^{N}x(t_i\oplus\Delta).
\]

\begin{proposition}[Uniform finite marginals and independent shift batches]
\label{prop:rqmc-finite-unbiased}
For every deterministic point set, $0\leq B(\Delta)\leq\rho$ and
$\mathbb E B(\Delta)=\mu$. If the shift vectors $\Delta_1,\ldots,\Delta_R$
are independent and each exactly uniform on $\mathcal G$, their batch means
$B(\Delta_1),\ldots,B(\Delta_R)$ are independent and identically distributed.
These statements do not imply independence among the points within a batch
or a variance reduction proportional to $N^{-1}$.
\end{proposition}

\begin{proof}
For every fixed $t\in\mathcal G$, the map $\Delta\mapsto t\oplus\Delta$
is its own inverse and therefore a bijection of the finite group. Consequently,
for every $u\in\mathcal G$,
$\mathbb P(t\oplus\Delta=u)=|\mathcal G|^{-1}$. Thus
$\mathbb E x(t\oplus\Delta)=\mu$. Linearity of expectation gives
$\mathbb E B(\Delta)=N^{-1}\sum_i\mu=\mu$, without using pointwise
independence. The range follows by averaging values in $[0,\rho]$.
Finally, each batch mean is the same deterministic function of its own shift
vector. Applying deterministic functions separately to independent identically
distributed vectors preserves independence and their common distribution.
\end{proof}

The engine forms its digital points as integers, truncating the retained
32-bit Joe--Kuo words to the relevant $b_j$ leading bits. It then draws an
independent $b_j$-bit shift for every coordinate of each replicate. It applies
exclusive-or directly to these finite words. No conversion to a floating-point
uniform variate enters the financial graph. The confidence argument above
uses only the finite-group bijection; it does not infer accuracy from the
quality or discrepancy of the chosen direction numbers. I retain the
author-distributed direction file, algorithm and BSD licence described in
the source ledger. My actual PCG64DXSM stream is a reproducible numerical
implementation of the declared ideal-shift model, and is not a proof of
that model's randomness assumptions.

\begin{proposition}[A range and moment prior for a shifted batch]
\label{prop:rqmc-prior-variance}
Suppose that the finite target has proved bounds
$\mathbb E x(U)\leq u$ and $\mathbb E x(U)^2\leq v$, where $U$ is uniform
on $\mathcal G$, $0\leq u\leq\rho$ and $0\leq v\leq\rho^2$.
For $\rho>0$, define
\[
 m_* = \min\{u,\rho/2,v/\rho\},\qquad
 V_* = m_*(\rho-m_*).
\]
Then $\operatorname{Var}(B)\leq V_*$. For $\rho=0$, $B=0$ and its variance
is zero. This bound requires no assumption on the within-batch covariance.
\end{proposition}

\begin{proof}
Convexity of the square gives, for every shift,
\[
 B(\Delta)^2\leq\frac1N\sum_i x(t_i\oplus\Delta)^2.
\]
The preceding proposition therefore yields $\mathbb E B^2\leq v$.
The pointwise inequality $B^2\leq\rho B$ gives a second bound
$\mathbb E B^2\leq\rho\mu$. Also $0\leq\mu\leq u$. Consequently,
\[
 \operatorname{Var}(B)\leq\min\{v,\rho\mu\}-\mu^2.
\]
To maximise this upper bound on $\mu\in[0,u]$, split the interval at
$v/\rho$. On $[0,v/\rho]$ the expression is
$\rho\mu-\mu^2$, whose derivative is $\rho-2\mu$. It increases until
$\rho/2$ and decreases afterwards. On $[v/\rho,u]$, when that interval is
nonempty, the expression is $v-\mu^2$ and is nonincreasing. The maximum is
therefore attained at $\mu=m_*$ and equals $V_*$. The degenerate case is
immediate from $0\leq B\leq0$.
\end{proof}

I use this deterministic prior only to select the graph and to cap the
confidence bound on its batch variance. My prospective selection count is
\[
 R_{\mathrm{plan}}=
 \max\left\{1,
 \left\lceil
 \frac{(2V_*+2\rho\epsilon/3)\log(2/\alpha_{\mathrm{event}})}
      {\epsilon^2}
 \right\rceil\right\},
 \qquad \epsilon=\frac{0.005}{H}.
\]
I minimise $64R_{\mathrm{plan}}$ times the retained conservative scalar-lane
operation count over both fully conditioned methods and all their inherited
prefixes. This count is a prospective graph-selection objective. It is not
the production stopping rule, a measured runtime, a globally optimal
classical cost, or a claimed $N^{-1}$ variance rate. I retain all candidates
and the deterministic tie-breaking rule before generating the production
shifts. The production confidence construction uses the prescribed
checkpoint allocation in the separate confidence proposition.

\subsection{Exact batch statistics and monetary reconstruction}
\label{sec:rqmc-integer-reconstruction}

In the retained experiment $N=64$ and $q\leq30$. For replicate $j$ I compute
\[
 T_j=\sum_{i=1}^{64}a(t_i\oplus\Delta_j),\qquad
 D=64\,2^q,\qquad B_j=T_j/D.
\]
Each summand is at most $2^q$, so $T_j\leq2^{36}$ and the 64-word unsigned
sum is exact. All cross-replicate totals and square totals are assembled
as arbitrary-precision integers. For each finite checkpoint I retain
$S=\sum_jT_j$, $S_2=\sum_jT_j^2$ and
$Q=\sum_{k=1}^{\lfloor R/2\rfloor}(T_{2k-1}-T_{2k})^2$. They determine
\[
 \overline B=\frac{S}{RD},\qquad
 \widehat\sigma_{\mathrm{pair}}^2=
 \frac{Q}{2\lfloor R/2\rfloor D^2},\qquad
 \widehat\sigma_{\mathrm{sample}}^2=
 \frac{RS_2-S^2}{R(R-1)D^2}.
\]
The ordinary sample variance is retained as a diagnostic. It does not
substitute for the proved pair-based confidence bound.

The square reduction uses $T=\ell+2^{16}h$, where
$0\leq\ell<2^{16}$ and $0\leq h\leq2^{20}$. In a reduction containing
at most $2^{20}$ words, $\sum\ell^2<2^{52}$,
$\sum\ell h<2^{56}$ and $\sum h^2\leq2^{60}$.
Each unsigned 64-bit reduction is therefore safe. I reconstruct
\[
 \sum T^2=\sum\ell^2+2^{17}\sum\ell h+2^{32}\sum h^2
\]
using Python integers. The same argument applies to the unsigned absolute
paired differences, which are at most $2^{36}$. I retain the integer batch
sums, shift-stream hash, financial-output hash and predetermined scalar
replay points. A separate interpreter replays every retained financial
output, every batch sum and every checkpoint moment.

If the confidence construction returns $\mu\in[\ell_\mu,h_\mu]$, and the
retained geometric correction centre is $c$ with total correction,
arithmetic and distribution allowance $\beta$, the undiscounted payoff
is enclosed by
\[
 [c+H\ell_\mu-\beta,\ c+Hh_\mu+\beta]\cap[0,H].
\]
I apply the strictly positive discount factor to this interval and report
an enclosed midpoint with its rounding allowance. For these frozen contracts,
$\beta\leq0.005$ and $r=0.03$, $T=1$. If the finite-mean confidence
half-width is at most $0.005/H$, the price interval has half-width at most
$e^{-0.03}(0.005+0.005)<0.01$, apart from the explicitly retained numerical
enclosure increments, which are smaller than the strict slack. If the
confidence intersection is empty, I retain that outcome as an explicit
confidence failure and produce no reconstructed price interval. If the
prespecified maximum is reached with a wider nonempty interval, I retain
the interval and record that the requested statistical certificate was
not obtained.
